
Existing hallucination detection methods are generally coupled to LLM generation at inference time. SelfCheckGPT \citep{Manakul2023SelfCheckGPT} draws multiple stochastic samples and measures consistency; verbalized confidence methods prompt LLMs to self-assess uncertainty; probing classifiers require white-box access to internal activations. Each of these creates a structural dependency: the LLM must be run at detection time, at a cost proportional to LLM inference.

An alternative question is whether the signal required for hallucination detection is already present in a single forward pass of a frozen, off-the-shelf NLI classifier applied to a (context, response) pair that already exists in a log. Post-hoc approaches using NLI models --- notably SummaC \citep{Laban2022SummaC} and TRUE \citep{Honovich2022TRUE} --- have demonstrated that DeBERTa-scale NLI classifiers can measure factual inconsistency between text pairs without generation. However, to our knowledge, no prior work has: (1) established AUROC baselines for frozen NLI applied directly to HaluEval across all three task types under a single unified evaluation; (2) characterized which NLI scoring strategy produces the strongest hallucination detection signal; or (3) provided a principled account of why NLI-based detection fails on certain task types while succeeding on others.

This paper presents results for ExtrospectiveNLI, which applies \texttt{cross-encoder/nli-deberta-v3-large} to HaluEval (context, response) pairs using net-contradiction scoring (P(contradiction) $-$ P(entailment)) with sentence-level max aggregation for dialogue and QA, and full-document response-level scoring for summarization. Three research questions are addressed:

\begin{itemize}
\item \textbf{RQ1 (Existence)}: Does frozen NLI achieve above-chance AUROC on HaluEval commission tasks?
\item \textbf{RQ2 (Mechanism)}: Is the score separation between hallucinated and non-hallucinated responses structural, or distributional artifact?
\item \textbf{RQ3 (Comparison)}: How does generation-free NLI compare to generation-based SelfCheckGPT-NLI?
\item \textbf{RQ4 (Boundary)}: Does the approach fail on omission-type tasks, and can this failure be explained structurally?
\end{itemize}

\textbf{Summary of findings.} On dialogue, AUROC = 0.709 (DeLong p $\approx$ 0, Cohen's d = 0.714). On QA, AUROC = 0.644 (DeLong p = 1.$29\times 10^{-282}$, Cohen's d = 0.779). On summarization, AUROC = 0.530 (DeLong p = 2.$02\times 10^{-13}$, Cohen's d = 0.220). The two commission tasks pass the pre-specified gate (AUROC > 0.55, DeLong p < 0.05); summarization fails. Wilcoxon rank-sum tests confirm score separation on all three tasks. The summarization failure is attributed to the mismatch between NLI contradiction detection and omission-type hallucinations. Design-choice ablations were not executed; the reported results characterize the full configuration but cannot attribute performance to individual components.
